\begin{proof}[Proof of \cref{obj:thm:margin-localization}]
\begin{enumerate}
\item \emph{The constant.}  Fix
\[
C:=|C_m|+2+u_0^{-\alpha}.
\]
\Cref{obj:ass:margin} gives \(0\le\alpha\), \(C_m>0\), and \(u_0>0\), so every summand is nonnegative and \(C>0\); moreover
\[
C\ge C_m+1
\qquad\text{and}\qquad
C\ge u_0^{-\alpha}.
\]
This constant depends only on \(C_m\), \(u_0\), and \(\alpha\).  It remains to prove the displayed inequality for an arbitrary well-formed observed law \(P\), policy class \(\Pi\) whose members are measurable, and every \(\pi\in\Pi\). % lean: margin_localization

\item \emph{Two elementary facts.}  Let \(D_\pi=\{x:\pi(x)\ne\pi^\star_P(x)\}\) as in \cref{obj:def:disagreement}.  By \cref{obj:thm:welfare-identity},
\[
R_P(\pi)=\int_{\mathcal X} |\tau_P(x)|\,\mathbf 1_{D_\pi}(x)\,dP_X(x).
\]
The integrand is nonnegative, so \(R_P(\pi)\ge0\).  The standing well-formed-law conditions give \(\tau_P=\mu_1-\mu_0\), and \cref{obj:ass:bounded-outcome} gives \(|\mu_0|\le1\) and \(|\mu_1|\le1\); hence \(|\tau_P|\le2\), which supplies the integrability used in the identity.  Also \(P_X(D_\pi)\le1\), since well-formedness makes \(P_X\) a probability measure. % lean: regret_eq_disagreement_integral

\item \emph{The basic decomposition.}  For every \(u\in(0,u_0]\),
\[
P_X(D_\pi)\le C_m u^\alpha+\frac{R_P(\pi)}{u}.
\]
Indeed, write
\[
Z=\{x:\tau_P(x)=0,\ x\in D_\pi\},\qquad
S_u=\{x:0<|\tau_P(x)|\le u\},\qquad
B_u=D_\pi\cap\{x:u<|\tau_P(x)|\}.
\]
Every \(x\in D_\pi\) lies in one of these three sets, according to whether \(\tau_P(x)=0\), \(0<|\tau_P(x)|\le u\), or \(|\tau_P(x)|>u\); hence
\[
D_\pi\subseteq Z\cup S_u\cup B_u .
\]
\Cref{obj:ass:zero-effect} gives \(P_X(Z)=0\): in its first alternative \(P_X\{\tau_P=0\}=0\), so its subset \(Z\) is null; in its second alternative, applied to \(\pi\in\Pi\), the set \(Z\) of zero-contrast disagreements is null by hypothesis.  \Cref{obj:ass:margin}, applied at the admissible window \(u\le u_0\), gives
\[
P_X(S_u)\le C_m u^\alpha .
\]
Finally, on \(B_u\) the integrand in the regret identity is at least \(u\), and the integrand is nonnegative, so restricting the integral to \(B_u\) gives
\[
u\,P_X(B_u)\le \int_{B_u}|\tau_P(x)|\,\mathbf 1_{D_\pi}(x)\,dP_X(x)\le R_P(\pi),
\qquad\text{hence}\qquad
P_X(B_u)\le \frac{R_P(\pi)}u .
\]
Monotonicity and subadditivity of \(P_X\) along the displayed inclusion now give the asserted decomposition.  The required measurability comes from the well-formed-law measurability of \(\tau_P\) together with the policy-measurability part of \cref{obj:ass:policy-class}; boundedness and finiteness are supplied by the well-formed-law identity \(\tau_P=\mu_1-\mu_0\), \cref{obj:ass:bounded-outcome}, and the probability property of \(P_X\). % lean: disagreement_measure_le_margin_plus_regret_over_u

\item \emph{The case \(\alpha=0\).}  Then \(\alpha/(1+\alpha)=0\) and \(R_P(\pi)^0=1\), while \(u_0^{-0}=1\) gives \(C=C_m+3\ge1\).  Hence
\[
P_X(D_\pi)\le1\le C=C\,R_P(\pi)^{\alpha/(1+\alpha)} .
\]
% lean: margin_localization

\item \emph{The case \(\alpha>0\), \(R_P(\pi)>0\), small window.}  Set
\[
u=R_P(\pi)^{1/(1+\alpha)}>0,
\]
and suppose first \(u\le u_0\).  Then the two exponent identities
\[
u^\alpha=R_P(\pi)^{\alpha/(1+\alpha)},
\qquad
\frac{R_P(\pi)}u=R_P(\pi)^{1-1/(1+\alpha)}=R_P(\pi)^{\alpha/(1+\alpha)}
\]
hold, so Step 3 gives
\[
P_X(D_\pi)
\le C_m u^\alpha+\frac{R_P(\pi)}u
=(C_m+1)\,R_P(\pi)^{\alpha/(1+\alpha)}
\le C\,R_P(\pi)^{\alpha/(1+\alpha)},
\]
the last step by \(C_m+1\le C\) and \(R_P(\pi)^{\alpha/(1+\alpha)}\ge0\). % lean: margin_localization

\item \emph{The case \(\alpha>0\), \(R_P(\pi)>0\), large window.}  Suppose instead \(u>u_0\).  Since \(\alpha>0\), raising to the power \(\alpha\) is increasing on nonnegative arguments, so
\[
u_0^\alpha\le u^\alpha=R_P(\pi)^{\alpha/(1+\alpha)}.
\]
Multiplying by \(u_0^{-\alpha}>0\) and using \(u_0^{-\alpha}\le C\),
\[
1=u_0^{-\alpha}u_0^\alpha
\le u_0^{-\alpha}\,R_P(\pi)^{\alpha/(1+\alpha)}
\le C\,R_P(\pi)^{\alpha/(1+\alpha)} .
\]
Combining with \(P_X(D_\pi)\le1\) from Step 2 gives the desired bound in this subcase. % lean: margin_localization

\item \emph{The case \(\alpha>0\), \(R_P(\pi)=0\).}  Let \(\varepsilon>0\), and put
\[
u=\min\left\{u_0,\left(\frac{\varepsilon}{C_m+1}\right)^{1/\alpha}\right\}.
\]
Then \(0<u\le u_0\), and monotonicity of positive powers gives
\[
u^\alpha\le \frac{\varepsilon}{C_m+1}.
\]
Step 3 with \(R_P(\pi)=0\) therefore gives
\[
P_X(D_\pi)\le C_m u^\alpha+\frac 0u
=C_m u^\alpha
\le \frac{C_m}{C_m+1}\,\varepsilon
\le \varepsilon.
\]
Since \(\varepsilon>0\) was arbitrary and \(P_X(D_\pi)\ge0\), \(P_X(D_\pi)=0\).  Also \(\alpha/(1+\alpha)>0\), so
\[
R_P(\pi)^{\alpha/(1+\alpha)}=0^{\alpha/(1+\alpha)}=0,
\]
and therefore
\[
P_X(D_\pi)=0\le C\,R_P(\pi)^{\alpha/(1+\alpha)} .
\]
% lean: margin_localization
\end{enumerate}
\end{proof}
